% 演示视频讲稿（口语版·PPT＋结果＋报告＋代码多窗口）——编译方式：XeLaTeX
\documentclass[11pt]{ctexart}
\usepackage[a4paper, margin=2.3cm]{geometry}
\usepackage{booktabs}
\usepackage{amsmath, amssymb}
\usepackage{enumitem}
\setlist{nosep}
\usepackage{parskip}
\usepackage{xcolor}
\newcommand{\cue}[1]{\noindent\colorbox{gray!15}{\parbox{\dimexpr\linewidth-2\fboxsep}{\textbf{#1}}}\par\smallskip}

\title{\textbf{演示视频完整讲稿（口语版）}\\
\large Chooser Option 定价项目 · 10--15 分钟 · 对着 PPT＋实验结果＋报告＋代码讲}
\author{佟铸城（Zhucheng Tong）}
\date{2026 年 9 月}

\begin{document}
\maketitle

\section*{录制前准备（约 5 分钟）}
按导师要求，演示要在 PPT、实验结果、报告和代码之间切换。提前打开并排好五个窗口（用多桌面或 Cmd+Tab 切换）：
\begin{enumerate}
  \item \textbf{PPT}：\texttt{docs/week8/final\_presentation.pptx}，用 PowerPoint 或 Keynote 全屏放映（13 页）；
  \item \textbf{工具}：终端先跑 \texttt{streamlit run app/streamlit\_app.py}，浏览器窗口放大；
  \item \textbf{报告}：\texttt{docs/week8/final\_report\_cn.pdf}，提前翻到第 5 页（第 4 章的价格对比图在这页）；
  \item \textbf{代码}：编辑器打开 \texttt{src/models/chooser.py} 和 \texttt{src/metrics.py}；
  \item \textbf{录屏}：按 Cmd+Shift+5，选“录制整个屏幕”，在“选项”里勾上麦克风；开专注模式屏蔽通知。
\end{enumerate}
先不录音干跑一遍，熟悉切屏路线；正式录 2--3 条，取最好的一条。

\section*{讲稿（总时长约 12--13 分钟）}

\cue{[0:00--1:00 · PPT 第 1--2 页（标题＋项目一览）]}
老师好，我是佟铸城。这个视频汇报 chooser option 定价项目这八周做的事。简单说，一共四步：先把数据管道搭起来，然后复现并验证了 BSM 的 chooser 定价，接着用很严格的方式测试了机器学习能不能把定价做得更准，最后做出了一个定价工具。先报三个最重要的数字：第一，论文的定价我复现出来是 29.13 美元；第二，BSM 假设波动率不变，这个假设的代价大概是价格的 35\%；第三，是一个负结果——在没见过的数据上，机器学习没有跑赢基准方法，这个结论之前也跟您确认过了。

\cue{[1:00--2:30 · PPT 第 3 页（chooser 原理）]}
先花一分钟讲清楚我们在给什么东西定价。Chooser option 的特点是：买的时候不用定方向，等到约定的选择日再决定它是看涨还是看跌——适合那种“肯定会有大行情，但涨跌说不准”的场景，比如财报或者美联储议息之前。定价的关键一步是：选择日那天，拿 call 和 put 里更值钱的那个；用 put-call 平价关系变形之后，整个合约可以拆成两样东西——一个到期日在 $T_2$ 的 call，加上一个到期日在选择日、行权价折算过的 put。拆完就是一个买好放着不用动的组合，所以在 BSM 框架下能直接写出定价公式。直觉上，它比同时买一个 call 加一个 put（也就是 straddle）要便宜，但比只买一个 call 或者只买一个 put 要贵——我们算出来的数字也正好是这个大小关系。

\cue{[2:30--3:30 · 切到代码 chooser.py（拆分函数＝第 19--26 行；教材公式版＝第 29--42 行；蒙特卡洛＝第 55--92 行，讲“小技巧”时把光标停在 62--66 行的注释上）]}
对应到代码。这个文件就是定价模块：上面这个函数就是刚才说的拆分，一共不到十行；下面还有另一个函数，是照 Hull 教材的公式独立再写的一遍——两个函数互相对答案，差别只在小数点后十五位，纯粹是计算机浮点数的误差，可以认为完全一致。再往下是蒙特卡洛模拟，这里有个小技巧：模拟到选择日之后，被选中的那个期权的价值不用接着往下模拟，直接套 BSM 公式算出来——这样一半的模拟噪声直接消失了，同样的模拟次数，误差只有普通做法的一半左右。

\cue{[3:30--4:30 · PPT 第 4 页（数据）]}
数据这块简单过一下：用了三个来源互为备份，2018 到 2024 年的日度数据，一共整理出 18 个特征，分成五类；情绪指标是按您的建议从 VIX 算出来的，用“VIX 在过去一年里处于什么位置”这个百分位数，天然就落在 0 到 1 之间。数据清洗的每个决定都写在文档里，比如那些特别极端的日子，我们一个个核对过，都是疫情崩盘这类真实发生的行情，所以只做了标记、没有删掉。整条数据管道挂在 GitHub 上，每周五自动更新，到今天还在自己跑。

\cue{[4:30--6:30 · PPT 第 5--6 页（复现＋基准）→ 切到报告 PDF 第 5 页的价格对比图（蓝灰双线那张），讲完回 PPT 第 6 页]}
第三、四周是理论这一侧的重头。复现这部分，我们用三条互相独立的证据交叉验证：两套公式对答案、两种模拟方法对答案、价格大小关系符合“不能白赚钱”的原则，全都通过。另外还多做了一件事：参考论文在选择日用的规则是“股价高于行权价就选 call”，这其实是个偷懒的近似——我们算了它到底亏多少：利率接近零的时候几乎不亏，但利率涨到 10\% 的时候，亏损放大十倍。也就是说，一个近似在某个环境下没事，换个环境可能就有事。然后是基准测试，思路是“事后对答案”：同一份合约定两次价，一次用当时能看到的历史波动率，一次用后来一年真实发生的波动率，两个价的差距就完全是“波动率不变”这个假设造成的。（切到报告这张图）蓝线是当时定的价，灰线是事后对出来的答案——2020 年初，灰线先涨起来了，蓝线还没反应，说明 BSM 在危机来之前定便宜了、危机过后又定贵了。这张图就把假设失效的过程讲完了。量化下来：就算用最好的历史估法，平均误差还是有价格的 35\%；而且短窗口的估法和长窗口的估法，出错的时段正好相反，没有哪一种在所有行情下都好用——这就是让机器学习上场的理由。

\cue{[6:30--7:40 · PPT 第 7--8 页（规矩与设计＋验证集）]}
机器学习部分，先讲规矩和设计。规矩很严：训练、验证、测试按时间先后切开；训练数据里会“偷看到未来”的部分全部删掉；测试集从第五周就封起来，只在最后打开一次。设计上走了个弯路：一开始让模型直接预测波动率是多少，七个模型全军覆没——训练数据是疫情那段大波动，验证数据是平静期，学到的东西过时了。后来改成只预测“相对当前波动率，该往上调还是往下调多少”：模型什么都不会时输出零，零就等于基准本身，怎么都不会差太多。（翻到第 8 页）这页是验证集的成绩单：图里唯一的绿色就是改造后的模型，赢了基准 43\%——但注意，这只是验证集。

\cue{[7:40--9:00 · PPT 第 9--10 页（测试集＋SHAP）]}
测试集打开之后，没有一个模型赢过基准：基准误差 8.51，最好的机器学习 14.75；验证集冠军从 2.8 掉到 17，是教科书里“对验证集过拟合”的活例子。原因三个：训练数据只经历了一轮半的市场起伏；测试期的答案里有 2025 年 4 月的大跌，谁也料不到；第五周那次胜利事后看有运气成分。反复检查过六遍，不是代码写错。（翻到第 10 页）SHAP 分析回答“模型到底学了什么”：最重要的信号是过去一年的波动率——它越低，模型越预测未来会涨回去，也就是均值回归。方向学对了，幅度定不准，所以问题不在算法，在数据里可学的完整周期太少——这跟您说的“单纯用时序模型硬拟合，不如定价模型推导得准”完全一致。

\cue{[9:00--11:30 · 切到浏览器 localhost:8501 的工具，页签顺序：定价（拖左栏行权价/期限滑块）→ 敏感性与压力（指情景表前两行）→ 趋势与成绩单（各停两三秒）→ 指一下侧栏“拉取最新数据”按钮（不用真点）→ 回 PPT 第 11 页（极端情况小结）]}
基于这个结论，工具是这么设计的。（定价页签）一个页面给两个价：主价格用全程表现最稳的方法算，也就是基准方法；旁边这个机器学习修正价只做参考，标着“在没见过的数据上未证明更准”的提醒。价格旁边的误差范围也不是随手拍的：市场平静的时候用平静时期的历史误差，恐慌的时候用恐慌时期的。拖一下行权价和期限，价格马上重算。（敏感性页签）极端情况测试：波动率涨一半，价格涨 49\%；但利率加两个百分点，价格只动 1.5\%。原因是这个组合里，call 和 put 对波动率的反应方向相同、会叠加，对利率的反应方向相反、互相抵消——所以它基本上就是一个“赌波动大小”的工具，对利率不敏感。（趋势和成绩单页签一句带过）第四周和第六周所有的评估表都存档在成绩单页签里；侧边栏这个按钮可以一键拉当天最新的数据。工程上还有个细节：界面代码里一行定价公式都没有，所有计算都在后台的服务模块里，还专门写了一条检查，防止界面里偷偷长出第二套公式。（回到 PPT 第 11 页）极端情况的完整数字在这一页，股价上涨或下跌 15\% 价格都会升高——这是它“不用选方向”的体现。

\cue{[11:30--13:00 · 切到代码 metrics.py（权重字典＝第 35--40 行；综合分函数＝第 83 行，一瞥即走）→ PPT 第 12 页（分工结论）]}
（切到 metrics.py）评估打分也是独立模块：八个指标，每个都除以基准方法的对应值再加权平均——好处是读数特别直观，综合分小于 1 就是赢了基准。整个八周，所有模型用的都是这同一把尺子。（回到 PPT）最后回答您提的 1+1 的问题。八周下来我的结论是：这两样东西该分工，不该比赛。定价模型负责搭框架——保证价格不出现自相矛盾、告诉你怎么对冲、错了能查出错在哪；机器学习负责框架照顾不到的零头。具体有五条路可以走：第一条“只学修正量”我们已经做了；第二条是给机器学习的修正乘一个折中系数，在“完全信基准”和“完全信模型”之间调节；第三条是让机器学习只做它擅长的事——判断现在市场是平静还是恐慌，然后由定价模型根据状态换参数；后面两条，一个是拿个股期权的隐含波动率来建模，一个是让机器学习给复杂定价模型当加速器，这两条需要新的数据，我理解是 enhancement 阶段最值得做的方向。

\cue{[13:00--13:40 · PPT 第 12--13 页（收尾＋谢谢）]}
总结一下：八周的四个节点都完成了。“最优定价方案”这个问题的答案是：在现在的数据条件下，用基准方法定价最靠谱；机器学习的框架已经搭好放在那里，等数据条件成熟就能接上。所有结果都固定了随机种子，一条命令就能从头复现，代码和每周的检查脚本都在仓库里。谢谢老师。

\section*{录后检查清单}
\begin{itemize}
  \item[$\square$] 总时长 10--15 分钟（正常语速约 12--13 分钟，超时可压缩 3:30 数据段）
  \item[$\square$] 回放开头 30 秒确认有声音；无通知弹窗、无个人信息入镜
  \item[$\square$] 文件命名 \texttt{chooser\_demo\_佟铸城.mov}；过大再找我压缩
\end{itemize}
\end{document}
